\documentclass[11pt, twoside, a4paper]{book}

% ===== ENCODING & FONTS =====
\usepackage[utf8]{inputenc}
\usepackage[T1]{fontenc}
\usepackage{libertine} % A beautiful classical serif font suitable for Latin texts
\usepackage{microtype} % Improves typography and spacing

% ===== PAGE GEOMETRY =====
\usepackage{geometry}
\geometry{
    a4paper,
    inner=3cm, % Margin for binding
    outer=2.5cm,
    top=3cm,
    bottom=3cm
}

% ===== HEADERS & FOOTERS =====
\usepackage{fancyhdr}
\pagestyle{fancy}
\fancyhead{} % Clear all headers
\fancyhead[RO,LE]{\thepage}
\fancyhead[CO]{\textit{Res Gestae Divi Augusti}}
\fancyhead[CE]{\leftmark}
\fancyfoot{} % Clear all footers
\renewcommand{\headrulewidth}{0.4pt}

% ===== CUSTOM MACROS FOR EPIGRAPHY =====

% 1. Original Stone Line Numbers (Small bold superscripts)
\newcommand{\linenum}[1]{\textsuperscript{\textbf{#1}}\,}

% 2. Large Chapter/Division Numbers (The prominent '1', '2' blocks)
\newcommand{\rgchapter}[1]{%
    \par\vspace{1.5em}\noindent
    {\Huge\bfseries #1}\quad\ignorespaces
}

% 3. Emendations/Restorations (Black & White alternative)
% In classical philology and epigraphy (Leiden Conventions), 
% text restored by an editor is placed within square brackets. 
% We use bold brackets to make them distinct without relying on color.
\newcommand{\emend}[1]{\textbf{[}#1\textbf{]}}

% Alternative: If you prefer italics instead of brackets, uncomment the line below:
% \renewcommand{\emend}[1]{\textit{#1}}

% 4. Section Symbol / Paragraph Breaks
\newcommand{\sect}{\S\ }

% ===== ADDITIONAL FORMATTING =====
\usepackage{setspace}
\onehalfspacing % Improves readability for classical commentaries
\usepackage{hyperref}
\hypersetup{
    colorlinks=false,
    pdfborder={0 0 0}
}

\begin{document}

% ===== FRONTMATTER =====
\frontmatter

\begin{titlepage}
    \centering
    \vspace*{4cm}
    {\Huge \textbf{Res Gestae Divi Augusti}\par}
    \vspace{1.5cm}
    {\Large An Epigraphic Transcription and Commentary\par}
    \vspace{4cm}
    {\Large \textit{Author Name}\par}
    \vfill
    {\large \today\par}
\end{titlepage}

\tableofcontents

\chapter{Preface}
This volume presents the text of the \textit{Res Gestae Divi Augusti}. 
To accommodate black-and-white printing, this edition avoids color-coding. 
Instead, following standard epigraphic practices (the Leiden Conventions), 
text that has been emended or restored from the damaged stones is enclosed 
in bold square brackets \emend{like this}.

Numbers in superscript (e.g., \linenum{5}) denote the original line breaks on the physical monument, while large numbers denote the traditional textual divisions.

\chapter{Introduction}
The \textit{Res Gestae Divi Augusti} (The Deeds of the Divine Augustus) is the funerary inscription of the first Roman emperor, Augustus, giving a first-person record of his life and accomplishments.

% ===== MAINMATTER =====
\mainmatter

\chapter{Text and Translation}

% The transcription below uses the macros defined in the preamble
% to match the elements from the reference image.

\rgchapter{1} 
Rérum gestárum díví Augusti, quibus orbem ter\emend{rarum} imperio populi Rom. subiecit, \sect et inpensarum, quas in rem publicam populumque \emend{Romanum} fecit, incísarum in duabus aheneís pílís, quae su\emend{nt} Romae positae, exemplar sub\emend{iectum}.

\vspace{1cm}

\linenum{1} Annós undéviginti natus exercitum priváto consilio et privatá impensá \linenum{2} comparávi, \sect per quem rem publicam \emend{do}minatione factionis oppressam \linenum{3} in libertátem vindicá\emend{vi. Quas ob res}\footnote{First footnote example corresponding to the blue superscript '1' in the original.} \emend{senatus de}cretís hono\emend{rificís in} \linenum{4} \emend{ordinem suum me adlegit C.~Pansa A.~Hirtio consulibus, consula}\linenum{5}\emend{rem locum sententiae dicendae simul dans,}\footnote{Second footnote example corresponding to the blue superscript '2'.} \emend{et imperium mihi dedit.} \sect \linenum{6} Rés publica \emend{ne quid detrimenti caperet, me pro praetore simul cum} \linenum{7} \emend{consulibus providere iussit.} \sect \emend{Populus autem eódem anno mé} \linenum{8} \textdegree\emend{consulem, cum cos.~uterque bello cecidisset, et trium virum reí publicae constituendae creavit.}

\vspace{1cm}

\rgchapter{2} 
\linenum{10} \textdegree Quí parentem meum \emend{interfecerunt, eó}s in exilium expulí iudiciís legi\linenum{11}timís ultus eórum \emend{facinus,} \sect e\emend{t posteá bellum inferentís reí publicae} \linenum{12} \emend{víci bis acie.}

% ===== BACKMATTER =====
\backmatter

\chapter{Commentary}
Detailed notes on the historical context, epigraphy, and the linguistic features of the text belong here.

\chapter{Bibliography}
\begin{thebibliography}{9}
    \bibitem{brunt}
    P.A. Brunt and J.M. Moore, \textit{Res Gestae Divi Augusti: The Achievements of the Divine Augustus}, Oxford University Press, 1967.
    
    \bibitem{cooley}
    Alison E. Cooley, \textit{Res Gestae Divi Augusti: Text, Translation, and Commentary}, Cambridge University Press, 2009.
\end{thebibliography}

\end{document}
